% =========================================================================
% CAPITOLO 11 — CONCLUSIONE
% =========================================================================

\chapter{Conclusione}
\label{ch:conclusion}

\noindent
{\rmfamily\itshape
\begin{tabular}[t]{@{}l@{\hspace{0.3cm}}p{12.5cm}@{}}
\textls[50]{\textbf{Trattati:}} &
\textls[50]{Sintesi · Risposta alla domanda di ricerca · Contributo metodologico ·
Trasferibilità · Lavori futuri e implicazioni di design}
\end{tabular}
}

\vspace{0.5cm}

Questa tesi è iniziata con una domanda apparentemente semplice: come si rende comprensibile un sistema quando non lo si può vedere? E con una seconda domanda che dà alla prima un metodo: in che misura una rappresentazione esperienziale alimentata dal machine learning può sostenere il lavoro della comprensione?

La Atlantic Meridional Overturning Circulation è servita da caso, ma la domanda non è chiaramente mai stata solo sull'AMOC. Riguardava la più ampia classe di sistemi --- invisibili, dinamici, che si dispiegano su scale temporali più lunghe di una vita umana --- il cui comportamento conta e la cui struttura resiste ai modi convenzionali in cui la disegniamo. Questo capitolo conclusivo riunisce l'arco, enuncia ciò che il lavoro contribuisce e segna le strade che apre.

% -------------------------------------------------------------------------
\section{Rispondere alla domanda di ricerca}
\label{sec:concl-answer}

Le tre fasi --- rilevare, tradurre, valutare --- hanno fatto tipi diversi di affermazioni (Capitolo~\ref{ch:rqs}) e si sono quindi chiuse in modo diverso; questo è parte della risposta.

La fase di \emph{rilevamento} (RQ1) ha prodotto affermazioni verificabili rispetto ai dati --- e hanno retto: la struttura verticale dell'AMOC si riduce a due modi che spiegano quasi il novantotto per cento della varianza, un autoencoder ha confermato che lo spazio ridotto è abbastanza piatto da fungere da sistema di coordinate (H2), e l'analisi di preallarme ha restituito un rigoroso risultato nullo (H3). Questi sono i risultati più solidi della tesi: affermazioni su un dataset.

La fase di \emph{traduzione} (RQ2) chiedeva solo che l'artefatto esistesse e funzionasse. \emph{Living Atlantic} lo fa: gira e mantiene le sue due condizioni allineate per provenienza, inquadramento e unità di misura attraverso un bundle di dati e una ricostruzione condivisi, differendo solo nella dimensione temporale oggetto di studio (Capitolo~\ref{ch:translate}).

La fase di \emph{valutazione} (RQ3) è stata più prudente, chiedendo come i partecipanti esprimessero la comprensione anziché se ottenessero punteggi più alti. La presentazione esperienziale non ha migliorato la comprensione in ogni caso; ha cambiato come i partecipanti comprendevano il sistema e quali aspetti sapevano identificare (Capitolo~\ref{ch:discussion}). Le due ipotesi ricevono quindi un sostegno direzionale anziché una prova: H-A (le rappresentazioni convenzionali possono lasciare strutture inaccessibili ai non specialisti) e H-B (la rappresentazione esperienziale può renderle più accessibili) trovano il loro sostegno più netto nella scena della ricorrenza, dove nessun partecipante della condizione statica ha identificato la struttura ricorrente mentre diversi hanno nominato ciò che mancava; il sostegno è più debole dove una struttura è già visibile in una figura immobile.

La risposta onesta non è che la visualizzazione esperienziale funzioni, ma che funziona \emph{condizionatamente}: la sua utilità dipende da ciò che va compreso, aiutando dove un pattern è difficile da vedere e offrendo meno dove è già visibile.

Come il Capitolo~\ref{ch:introduction} anticipava, il contributo principale è metodologico, non oceanografico: un framework trasferibile, qui evidenziato attraverso l'AMOC, per trasformare un sistema complesso e difficile da osservare in una rappresentazione testabile, con ciascuna fase tenuta al proprio standard di evidenza --- verifica, costruzione e confronto esplorativo.

Due ulteriori contributi gli stanno accanto. Empiricamente, lo studio ha fatto emergere un errore ricorrente: più della metà dei partecipanti ha letto l'altezza della curva cumulata della funzione di corrente come velocità locale, una lettura errata che la ricchezza aggiunta non ha ridotto e che era più rara tra i partecipanti più esperti --- inatteso, e da testare su larga scala. Materialmente, \emph{Living Atlantic} resta un oggetto di design riproducibile costruito sulla provenienza condivisa, utilizzabile in studi futuri anziché solo come dimostrazione.

Insieme, questi collocano la tesi come un ponte tra dati scientifici e comprensione umana anziché come un contributo all'oceanografia, allineata come esposto nella \S\ref{sec:contribution} con l'\textbf{SDG~13} (educazione e consapevolezza climatica), l'\textbf{SDG~4} e l'\textbf{SDG~14}.

% -------------------------------------------------------------------------
\section{Lavori futuri}
\label{sec:concl-future}

Il passo successivo più chiaro discende dai risultati. Il piccolo campione di convenienza e la differenza di alfabetizzazione visiva pregressa tra i gruppi limitano le conclusioni; uno studio più ampio, bilanciato su quell'esperienza pregressa e che controbilanci l'ordine delle scene (\S\ref{sec:eval-limitations}), testerebbe più chiaramente se la visualizzazione esperienziale aiuti per alcuni tipi di struttura più che per altri --- trattando le tre scene come test separati di un'unica domanda: quando la presentazione esperienziale migliora davvero la comprensione?

Anche l'errore suggerisce un miglioramento di design. Poiché l'altezza della curva è facilmente letta come velocità, una versione futura dovrebbe rendere visibile la sua natura cumulata --- per esempio un'animazione che costruisce l'integrale mentre il cursore scende lungo la colonna, mostrando accanto la pendenza locale come la quantità che porta direzione e intensità. Questo affronta direttamente l'errore più comune dello studio.

Più in generale, il framework potrebbe essere applicato oltre l'AMOC: molti sistemi importanti sono difficili da osservare direttamente, e testare lo stesso approccio su di essi mostrerebbe se si trasferisce.

% -------------------------------------------------------------------------
\section{Chiusura}
\label{sec:concl-closing}

La missione che ha aperto questo lavoro --- \emph{Wisdom Unlocks Awareness, and Awareness Unlocks Agency} --- era un'affermazione su una sequenza: che la comprensione trasforma la conoscenza nella capacità di agire. Questa tesi ha testato il primo anello e l'ha trovato condizionato. La comprensione può essere costruita, ma non rendendo un sistema meramente visibile: dipende da quale struttura è trasmessa, a chi, e a quale costo in attenzione. È un'affermazione più piccola di quella da cui sono partito, e più utile.

La tesi si proponeva di rendere intelligibile un sistema invisibile, e di chiedere se cambiare come esperiamo i dati cambi come li comprendiamo. Lascia quella domanda non risolta una volta per tutte, ma risposta con precisione: per alcune strutture, per alcuni osservatori, e a un costo percettivo attorno a cui bisogna progettare. L'invisibile non è diventato pienamente visibile. Ma è diventato, in alcuni punti e per ragioni che lo studio può nominare, un po' più intelligibile --- e ha mostrato esattamente dove il prossimo tentativo dovrebbe mirare.
